Clouds play a vital role in regulating Earth's radiative budget and hydrological cycle. They interact with solar and thermal radiation, affecting the energy balance and, consequently, the surface temperature \citep{twomeyInfluencePollutionShortwave1977, brenguierCloudMicrophysicalRadiative2003, zengModelingRadiativeEffect2021}. Additionally, clouds transport moisture to various regions, being directly associated with drought, which is intensifying in the Mediterranean Region \citep{rios-entenzaMoistureRecyclingIberian2014, hoerlingIncreasedFrequencyMediterranean2012a}. However, studying clouds is quite complex, and their radiative effects are in the spotlight of the atmospheric science community. This challenge arises from clouds' high spatial and temporal variability, as well as the variety of hydrometeor types inside them \citep{intergovernmentalpanelonclimatechangeipccEarthEnergyBudget2023}. The hydrometeor phase, shape and concentration can significantly alter their interaction with solar and thermal radiation \citep{yoshidaEffectsVerticalProfiles2005}, posing a significant challenge to estimate their radiative effect as suggested by \cite{XXX}.


The ground-based remote sensing technique is one of the state-of-the-art methods for measuring clouds. Particularly, the Cloud Doppler Radar (CDR) can perform long-term measurements of clouds, which is one alternative for improving the cloud knowledge. Several studies have been analyzing CDR data in Europe, but they are mostly located in Northern Europe \citep{kneifelMultiyearCloudPrecipitation2022, nomokonovaStatisticsCloudsTheir2019, buhlMeasuringIceLiquidwater2016, achtertPropertiesArcticLiquid2020}. In this context, the AGORA (Andalusian Global ObseRvatory of the Atmosphere) ACTRIS CCRES station has been operating a 94 GHz CDR since 2018, contributing to enhancing the lack of data on the Mediterranean basin, which is a key region for climate variability. This station is one of the southernmost stations that is part of the Cloudnet network \citep{illingworthCloudnet2007}. Notably, it stands as the only cloud remote sensing database available in Spain, making it an invaluable resource for regional climate studies. 
Satellite data would be an alternative to fill the gap of data, but they only provide a general  picture of cloud properties. They present low temporal resolution and information on low level clouds is affected by large uncertainties. Additionally, this data should be validated with the use of ground-based remote sensors such as lidars and CDRs \citep{protatImpactConditionalSampling2006, protatEvaluationCloudSatCALIPSO2010}.
Thus, ground-based observations are still necessary for improving the accuracy of cloud properties.

This study aims to perform a statistical analysis on single layer clouds for more than half a decade of observations at the AGORA station. First, clouds were classified (e.g. ice, liquid or mixed-phase clouds) using the cluster classification algorithm (CBA). This is a novel algorithm that classifies clouds according to the percentage of hydrometeors within each cluster. The clusters are defined using the hydrometeor classification product from Cloudnet.
Second, cloud macrophysics (e.g., cloud thickness, base and top height) and microphysics statistics are performed for each cloud category from CBA. Markedly, cloud statistics in this study are derived from a cluster-based approach, contrasting with previous studies that used the profile-based algorithm (PBA) \citep{pirloagaGroundBasedMeasurementsCloud2022,nomokonovaStatisticsCloudsTheir2019,kneifelMultiyearCloudPrecipitation2022}.
For instance, \cite{nomokonovaStatisticsCloudsTheir2019} used the PBA algorithm to characterize single-layer clouds at the Ny-Ålesund station, \cite{kneifelMultiyearCloudPrecipitation2022} employed it to investigate ice cloud variability in the German Alps, and \cite{pirloagaGroundBasedMeasurementsCloud2022} applied it to examine seasonal cloud variability in Bucharest. \cite{buhlMeasuringIceLiquidwater2016} also uses the PBA but selected intervals with more than 15 min with the same cloud category. On the other hand, the CBA considers the composition of the entire volume occupied by the cloud. \cite{roschkeDiscriminatingDrizzleRain2024} also used the CBA, but with a different classification scheme, separating clusters into warm and cold clouds for the Barbados site. 


The experimental site, datasets, instruments and products are described in Section \ref{sec:site_instruments_products}. The proposed new approach on cloud classification algorithm and the cloud assessment are presented in Section \ref{sec:classification_algorithm}. Section \ref{sec:results} analyzes the annual, seasonal, and daily cycle of cloud types, as well as the base and top height, and thickness. In addition, this section discusses the liquid and ice properties variability for different clouds. Finally, Section \ref{sec:conlusion} summarizes the main findings and discusses their importance to improve further cloud statistics analysis.
